\begin{abstract}
\noindent\textbf{Purpose:} Adolescent idiopathic scoliosis (AIS) clusters at peak height velocity, in the thoracic spine, and overwhelmingly in girls (8:1), yet no single framework unifies its timing, location, and sex predominance. We test whether AIS can be modelled as a failure to recover spinal geometry during growth --- when growth-driven remodeling outpaces active geometric restoration.

\noindent\textbf{Methods:} We model the developing spine as a biomechanical control system: a Cosserat rod whose segment-wise stiffness is set by developmental HOX-domain patterning and stabilised against gravity by proprioceptive feedback. Four research strands: (i) cross-species allometric analysis across 10 vertebrates (mouse to elephant) using $\mathcal{B}_g = EI/MgL^2$; (ii) AlphaFold-derived structural metrics (elongation, hinge-density per 100 residues; length-matched controls) for 61 candidate proteins; (iii) a two-timescale recovery-ratchet model of curve setting through the growth spurt (geometric-restoration rate $k_r$ versus growth-remodeling rate $k_g$); (iv) calibration against a simulated cohort ($N=1,000$) and bracing simulation compared against the Lonstein \& Carlson progression index and the BrAIST clinical trial.

\noindent\textbf{Results:} Cross-species, $\mathcal{B}_g$ declines monotonically with body mass (exponent $-0.282$, SE $0.058$, $R^2=0.744$), but we report as a negative result that it does \emph{not} isolate the human condition: the human adult ($\mathcal{B}_g = 0.0101$) is indistinguishable from the rabbit and sits above cat, dog, horse and elephant, so the case for obligate active maintenance rests on posture and axial load path rather than on a $\mathcal{B}_g$ threshold. A curated subset of extended-filament/channel mechanosensors (Vimentin, Lamin~A/C, PIEZO2, SCRIB, PTK7) is 72\% more structurally elongated than metabolic regulators ($p=0.011$, Cohen's $d=1.19$, $n=23$; signal dilutes with broader gene sets) and has 70\% fewer hinge regions per residue ($p=0.0025$, $n=61$) --- consistent with cooperative load-sensing. In the recovery-ratchet simulation, the growth-remodeling rate peaks at peak height velocity and overtakes the geometric-restoration rate, so transient curvature is progressively cemented into permanent set --- predicting that curve flexibility (reducibility) declines as the curve progresses; over the peak-height-velocity interval the model-derived demand-to-supply ratio rises by ${\sim}29\%$, a prefactor-free relative-growth prediction rather than a calibrated threshold. The effective Bio-Gravitational parameter $B_g^{\mathrm{eff}}$ tracks the Lonstein \& Carlson Progression Factor across the simulated cohort ($r = -0.483$); both quantities are deterministic functions of the same four simulated demographic variables, so this is a structural consistency check rather than independent validation, and the nominal $p$-value is not interpretable as evidence. Simulated bracing (modeled as 37\% gravitational load-sharing) reproduces the BrAIST success rates (72\% braced vs 48\% observation), which were themselves the grid-search targets used to calibrate the three free parameters; the Risser-stratified subgroup rates were \emph{not} calibration targets and replicate in direction but not magnitude (model 70.3\%/74.4\% vs trial 64\%/78\%). In a deterministic sweep, the critical spine length $L_{\mathrm{crit}}$ \emph{declines} with structural anisotropy $\mathcal{A}$ ($r=-0.88$), identifying reduction of $\mathcal{A}$ as the therapeutic direction; $20$ of $50$ sweep points lie on the search bounds, so this is reported as a direction rather than a calibrated magnitude.

\noindent\textbf{Conclusions:} The framework reframes AIS as a predictable failure of an active biomechanical control system, integrating developmental patterning, proprioceptive feedback, and tissue mechanics. The results show that the Bio-Gravitational parameter serves as a robust biomechanical proxy for clinical progression and brace-treatment efficacy. Three hypothesis-generating predictions follow, presented \emph{not} as clinical recommendations: (i) curve flexibility (reducibility on supine and bending films) should decline through the growth spurt in adolescents who progress, and low reducibility at presentation should predict progression independent of skeletal maturity; (ii) the model predicts that mitochondrial NAD$^+$ availability modulates simulated curve progression --- \textbf{this does not endorse nutraceutical use outside an IRB-approved trial}; (iii) ciliary length and cellular dithering should increase in scoliotic osteoblasts before macroscopic onset.
\end{abstract}

